% ======================================================================
% 4. PASS-THE-HASH IN DEPTH  (budget: 3.5 pages)
% Job of the chapter: the reader can reproduce the response computation on
% paper, knows why the wire evidence looks innocent, and knows which
% protocol choice changes the outcome.
% ======================================================================
\section{Pass-the-Hash in Depth}
\label{sec:pth}

\subsection{What the attacker actually holds}
% CLAIM (0.4): sixteen bytes, and their limits.
% Beats: no salt, no key history, no plaintext; where it comes from in
% practice (LSASS, SAM/SECURITY from a backup or shadow copy, a DC dump);
% what it does NOT buy (plaintext, AES Kerberos keys, reuse elsewhere).

\subsection{Computing a valid response without the password}
\subsubsection{The derivation, step by step}
% CLAIM (0.5): the response is HMAC-MD5 keyed with the NT hash over server
% challenge || client challenge || timestamp || target info. Walk it once,
% explicitly, and stop -- the rest of the section is consequences.
\subsubsection{Why this is not replay}
% CLAIM (0.3): fresh timestamp and client challenge per attempt, and the
% binding to the server's challenge, mean a captured response cannot be
% reused against another host. The proof lives in how the bytes were
% produced, not in the bytes.
\subsubsection{Two ways to pass it}
% CLAIM (0.3): speak the protocol from the attacker's host (network logon),
% or write the hash into LSASS and let Windows do it. Different logon types,
% different forensic footprints -- feeds chapters 7 and 8.
% Figure: fig-locality -- what crosses the wire vs what never does.

\subsection{Consequences for signing and sealing}
% CLAIM (0.5): the session key derives from the same hash, so an attacker
% holding the hash can also sign and seal. SMB signing therefore defeats
% relay, not PtH. This is the single most frequently botched claim in
% student write-ups; make it with the derivation, not with an assertion.

\subsection{Protocol coverage}
\subsubsection{SMB}
% CLAIM (0.2): named pipes, ADMIN$/IPC$, service creation or scheduled task;
% what privilege the account needs and what the target records.
\subsubsection{WinRM}
% CLAIM (0.2): no share dependency; often permitted by firewall policy where
% SMB is blocked. Frequently the more reliable path in hardened estates.
\subsubsection{WMI and DCOM}
% CLAIM (0.2): RPC over 135 plus a dynamic range; no file share required.
\subsubsection{RDP with Restricted Admin}
% CLAIM (0.2): a special case worth its own paragraph -- the logon is
% network-type and the hash never reaches the target, it is consumed locally.
\subsubsection{Secondary carriers}
% CLAIM (0.2): MSSQL, LDAP, HTTP/Negotiate. Short; the point is that the
% technique is protocol-agnostic and therefore cannot be fixed per-protocol.

\subsection{Tooling landscape}
% CLAIM (0.4): a table mapping capability to protocol, not a tool review.
% Note our position: orchestrate maintained implementations, own the glue and
% the evidence pipeline. No reimplementation of crypto primitives.
